\textit{Bearbeitet von Masoud Fahim}\\

Die Aufgabenstellung sieht zwei Regler vor: einen Positionsregler für die Öffnungsweite
und einen Drehmomentregler für die Greifkraft. Der Positionsregler wurde wie vorgesehen
entworfen und umgesetzt. Für die Greifkraft wurde zunächst ein geschlossener Kraftkreis
über die StallGuard-Funktion des Treibers entworfen und simuliert; die Messungen am
Prüfstand zeigten jedoch, dass das Messglied die auftretende Last nicht auflöst. Die
Kraftaufgabe wird deshalb über den Positionsregler und die bekannte Federkennlinie
gesteuert. Dieses Kapitel beschreibt beide Wege, weil erst der verworfene Entwurf die
gewählte Lösung begründet.

% +++++++++++++++++++++++++++++++++++++++++++++++++++++++++++++++++
% Regelstrecke und Reglerstruktur
% +++++++++++++++++++++++++++++++++++++++++++++++++++++++++++++++++
\section{Regelstrecke und Wahl der Reglerstruktur}
\label{sec:strecke}

Der Motortreiber TMC2209 wird nicht über STEP/DIR, sondern über das Register
\texttt{VACTUAL} angesteuert. In dieser Betriebsart gibt der Mikrocontroller eine
Sollgeschwindigkeit vor, und der Treiber erzeugt die zugehörigen Mikroschritte
eigenständig mit einem internen Taktgenerator. Die Stellgröße des Regelkreises ist damit
eine Geschwindigkeit und keine Position.

Aus dieser Eigenschaft folgt die gesamte weitere Auslegung. Die Position der Zahnstange
ergibt sich als Zeitintegral der kommandierten Geschwindigkeit, die Regelstrecke ist also
ein reiner Integrator:

\begin{equation}
    G(s) = \frac{X(s)}{V(s)} = \frac{1}{s}
    \label{eq:strecke}
\end{equation}

Mit einem Proportionalregler $K_P$ lautet der offene Kreis

\begin{equation}
    L(s) = K_P \cdot G(s) = \frac{K_P}{s} ,
    \label{eq:offenerkreis}
\end{equation}

woraus sich die Durchtrittsfrequenz unmittelbar ablesen lässt. Der Betrag von $L(s)$
erreicht bei $\omega = K_P$ den Wert eins, sodass

\begin{equation}
    \omega_c = K_P
    \label{eq:durchtritt}
\end{equation}

gilt. Die Reglerverstärkung entspricht damit direkt der angestrebten Bandbreite des
geschlossenen Kreises, was die Parametrierung auf eine einzige Entwurfsentscheidung
reduziert.

\subsection*{Begründung der reinen P-Struktur}

Die Aufgabenstellung nennt eine Auslegung als PI-, PID- oder Zustandsregelung. Umgesetzt
wurde ein reiner P-Regler. Diese Abweichung ist keine Vereinfachung, sondern folgt
zwingend aus der Struktur der Strecke.

Der Integrator, der bei einer PI-Regelung üblicherweise im Regler ergänzt wird, um
bleibende Regelabweichungen zu vermeiden, ist hier bereits in der Strecke enthalten. Die
stationäre Genauigkeit ist deshalb ohne zusätzlichen I-Anteil gegeben: Damit die Position
im Beharrungszustand konstant bleibt, muss die kommandierte Geschwindigkeit null sein.
Wegen $v = K_P \cdot e$ ist das gleichbedeutend mit einer Regelabweichung von null. Ein
reiner P-Regler arbeitet an dieser Strecke also stationär exakt.

Ein zusätzlicher I-Anteil würde dagegen einen zweiten Integrator in den offenen Kreis
einbringen. Der Phasengang läge damit bereits ohne weitere Verzögerungen bei
$-180^\circ$, und die Stabilität des Kreises hinge allein von der Wahl der
Reglernullstelle ab. Der Regler wäre schwieriger zu parametrieren, ohne einen Vorteil bei
der Genauigkeit zu bieten.

Als Nebeneffekt entfällt eine sonst notwendige Anti-Windup-Maßnahme. Bei großen
Sollwertsprüngen arbeitet der Regler über den größten Teil der Fahrstrecke in der
Stellgrößenbegrenzung. Ein I-Anteil würde in dieser Zeit einen erheblichen Fehler
aufsummieren und nach Erreichen der Zielposition ein Überschwingen verursachen. Da kein
Integralanteil vorhanden ist, tritt dieser Effekt nicht auf, und die Begrenzung kann ohne
Zusatzlogik implementiert werden.

% +++++++++++++++++++++++++++++++++++++++++++++++++++++++++++++++++
% Positionsregler
% +++++++++++++++++++++++++++++++++++++++++++++++++++++++++++++++++
\section{Positionsregler}
\label{sec:positionsregler}

\subsection*{Wahl der Bandbreite}

Nach Gleichung \ref{eq:durchtritt} entspricht die Reglerverstärkung unmittelbar der
Bandbreite des geschlossenen Kreises. Die Auslegung reduziert sich damit auf die Frage,
wie schnell der Kreis sein soll. Als Zielbandbreite wurden $f_c = 20$~Hz gewählt, woraus

\begin{equation}
    K_P = 2\pi f_c = 2\pi \cdot 20\,\mathrm{Hz} = 125{,}7\,\mathrm{s^{-1}}
    \label{eq:kp}
\end{equation}

folgt. Der geschlossene Kreis verhält sich damit im linearen Bereich wie ein
Verzögerungsglied erster Ordnung mit der Zeitkonstante

\begin{equation}
    \tau = \frac{1}{K_P} = 7{,}96\,\mathrm{ms} .
    \label{eq:tau}
\end{equation}

Der Reglertask läuft mit einer Abtastzeit von $T = 2$~ms, entsprechend 500~Hz. Das
Produkt $K_P \cdot T = 0{,}251$ liegt deutlich unter eins, sodass der Kreis pro Abtastschritt
nur einen Bruchteil der verbleibenden Abweichung abbaut und die Diskretisierung das
Verhalten nicht dominiert.

\subsection*{Stellgrößenbegrenzung}

Die Geschwindigkeit wird auf $v_{max} = 40$~mm/s begrenzt. Dieser Wert ergibt sich aus der
zulässigen Motordrehzahl und entspricht einem Registerwert von \texttt{VACTUAL}~$= 3351$.

Der Begrenzung nachgeschaltet ist eine Anstiegsbegrenzung mit $a_{max} = 5000$~mm/s².
Sie ist kein optionales Komfortmerkmal: Der TMC2209 besitzt keinen Rampengenerator und
setzt eine sprunghaft geänderte Sollgeschwindigkeit unmittelbar um. Ohne softwareseitige
Begrenzung würde der Antrieb bei jedem Sollwertsprung mit einem Geschwindigkeitssprung
anfahren, was zu Schrittverlusten führen kann. Die Anstiegsbegrenzung übernimmt damit
die Funktion, die bei anderen Treibern in der Hardware liegt.

Der Zahlenwert ist nach unten begrenzt durch den Regler selbst. Bei einem Sollwertsprung
der Höhe $\Delta x$ ändert der Regler seine Ausgangsgröße mit
$\dot{v} = K_P^2 \cdot \Delta x$, bei einem Sprung von 0,2~mm also mit rund
3160~mm/s². Eine Anstiegsbegrenzung unterhalb dieses Wertes würde den Regler ausbremsen
und den Kreis verlangsamen, statt ihn zu schützen.

\subsection*{Totzone und Grenzzyklus}

Ohne weitere Maßnahme kommt der Antrieb im Stillstand nicht zur Ruhe. Ursache ist das
Verhältnis der beiden Auflösungen im Kreis. Ein Mikroschritt entspricht an der Zahnstange

\begin{equation}
    \Delta x_{\mu} = \frac{\pi \cdot d}{n_{voll} \cdot n_{\mu}}
    = \frac{\pi \cdot 17\,\mathrm{mm}}{200 \cdot 16} = 0{,}01669\,\mathrm{mm} ,
    \label{eq:mikroschritt}
\end{equation}

während eine Stufe des Encoders

\begin{equation}
    \Delta x_{enc} = \frac{\pi \cdot d}{2^{12}}
    = \frac{53{,}407\,\mathrm{mm}}{4096} = 0{,}01304\,\mathrm{mm}
    \label{eq:encoderstufe}
\end{equation}

beträgt. Der kleinstmögliche Stellschritt des Aktors ist damit um den Faktor 1,28 größer
als die kleinste auflösbare Messgröße. Der Motor kann folglich nicht auf einer
Encoderstufe zum Stehen kommen, sondern springt über sie hinweg. Der Regler erkennt die
entstandene Abweichung, korrigiert zurück, springt erneut darüber und so fort. Das
Ergebnis ist ein Grenzzyklus, der sich als dauerhaftes Pendeln der Geschwindigkeit um den
Nullpunkt äußert und am realen Antrieb als Brummen im Stillstand hörbar wird.

Abhilfe schafft eine Totzone unmittelbar hinter dem Summierer, die Abweichungen unterhalb
einer Schwelle nicht mehr korrigiert. Rechnerisch genügen dafür zwei Encoderstufen, also
0,026~mm. Am realen Aufbau erwies sich dieser Wert als zu klein: Die Herleitung
berücksichtigt ausschließlich die Quantisierung, nicht jedoch das Rauschen des
Messsignals. Die Totzone wurde deshalb auf 0,05~mm erhöht, was etwa vier Encoderstufen
entspricht. Erst mit diesem Wert geht die Geschwindigkeit im Stillstand auf null.

Die Totzone führt eine bleibende Positionsabweichung von bis zu 0,05~mm ein. Dieser
Kompromiss ist bewusst gewählt und nicht regelungstechnisch, sondern durch die Hardware
begründet: Feiner zu regeln, als der Aktor stellen kann, führt zwangsläufig zum
Grenzzyklus. Gegenüber der geforderten Positioniergenauigkeit von $\pm 0{,}5$~mm bleibt
ein Sicherheitsfaktor von zehn.

\subsection*{Aufbereitung des Messsignals}

Vereinzelt traten auf dem über I$^2$C gelesenen Encoderwert einzelne stark abweichende
Messwerte auf. Da der Regler proportional auf die Abweichung reagiert, würde ein solcher
Ausreißer unmittelbar in einen Geschwindigkeitssprung umgesetzt. Dem Encoderwert ist
deshalb ein Medianfilter über drei Werte nachgeschaltet. Er unterdrückt einzelne Spikes
vollständig, ohne die Dynamik des Kreises nennenswert zu verzögern, da er im Gegensatz zu
einem Mittelwertfilter keine Glättung über mehrere gültige Messwerte vornimmt.

\subsection*{Stabilitätsgrenzen}

Der Entwurfswert von $K_P = 125{,}7\,\mathrm{s^{-1}}$ ist nach oben durch zwei Effekte
begrenzt, die im kontinuierlichen Modell nach Gleichung \ref{eq:strecke} nicht enthalten
sind.

Die Abtastung führt eine Grenze ein, oberhalb derer der diskrete Kreis instabil wird. Für
$K_P \cdot T > 1$ überschreitet der Kreis das Ziel innerhalb eines Takts, ab
$K_P \cdot T = 2$ läuft er auf. Bei $T = 2$~ms entspricht das einer Grenze von
$K_P = 1000\,\mathrm{s^{-1}}$.

Die zweite Grenze stammt aus der Motordynamik. Der Rotor folgt dem Statorfeld nicht starr,
sondern als Feder-Masse-System aus magnetischer Steifigkeit und Trägheit. Bei dessen
Eigenfrequenz $\omega_0$ trägt dieses Schwingungsglied zusätzlich $-90^\circ$ zur bereits
vorhandenen Phasendrehung des Integrators bei, sodass der Kreis dort auf der kritischen
Phasenlage liegt. Über die Stabilität entscheidet dann allein der Betrag, woraus sich

\begin{equation}
    K_{P,max} = 2 \zeta \omega_0
    \label{eq:resonanzgrenze}
\end{equation}

ergibt. Mit $\omega_0 = 628\,\mathrm{rad/s}$ liefert diese Bedingung je nach Dämpfung
sehr unterschiedliche Werte, wie Tabelle \ref{tab:stabilitaetsgrenzen} zeigt.

% ---------------------------------------------------------------
\begin{table}[htbp]
    \centering
    \caption{Grenzen für $K_P$ aus Abtastung und Motorresonanz}
    \label{tab:stabilitaetsgrenzen}
    \begin{tabular}{lcc}
        \hline
        Ursache & Grenze für $K_P$ & Reserve zum Entwurfswert \\
        \hline
        Abtastung, $T = 2$~ms          & 1000 & Faktor 8 \\
        Motorresonanz, $\zeta = 0{,}20$ & 251  & Faktor 2 \\
        Motorresonanz, $\zeta = 0{,}10$ & 126  & keine \\
        \hline
    \end{tabular}
\end{table}
% ---------------------------------------------------------------

Die schärfere der beiden Bedingungen ist damit nicht die Abtastung, sondern die Resonanz.
Bemerkenswert ist der Fall $\zeta = 0{,}1$: Die verbreitete Faustregel, die Bandbreite
höchstens auf ein Fünftel der Resonanzfrequenz zu legen, ist rechnerisch identisch mit
Gleichung \ref{eq:resonanzgrenze} für genau diese Dämpfung. Die Faustregel unterstellt
also stillschweigend einen Zahlenwert, der im vorliegenden Fall nicht bekannt war.

Die Dämpfung wurde nicht messtechnisch bestimmt. Für die Inbetriebnahme wurde deshalb ein
reduzierter Startwert von $K_P = 100\,\mathrm{s^{-1}}$ vorgesehen. Der anschließende
Betrieb bei $K_P = 125{,}7\,\mathrm{s^{-1}}$ verlief ohne erkennbare Schwingneigung, was
nach Gleichung \ref{eq:resonanzgrenze} auf eine Dämpfung oberhalb von 0,1 schließen lässt.
Die Auslegung ist damit nachträglich abgesichert, ohne dass $\zeta$ als Zahlenwert
vorliegt.

\subsection*{Struktur des Regelkreises}

Abbildung \ref{fig:positionsregelkreis} fasst die beschriebenen Elemente zusammen.

% ---------------------------------------------------------------
\begin{figure}[htbp]
\centering
\begin{tikzpicture}[node distance=5mm,
    blk/.append style={text width=14mm}]

  \node (soll)                        {};
  \node[sum, right=8mm of soll] (s)   {$\pm$};
  \node[blk, right=6mm of s]    (dz)  {Totzone\\[-1pt]\tiny 0,05\,mm};
  \node[blk, right=3mm of dz]   (kp)  {$K_P$\\[-1pt]\tiny 125,7\,s$^{-1}$};
  \node[blk, right=3mm of kp]   (sat) {Sättigung\\[-1pt]\tiny $\pm$40\,mm/s};
  \node[blk, right=3mm of sat]  (rl)  {Rate Limiter\\[-1pt]\tiny 5000\,mm/s$^2$};
  \node[blk, right=3mm of rl]   (tmc) {Aktor\\[-1pt]\tiny TMC2209};
  \node[blk, right=3mm of tmc]  (int) {Integrator\\[-1pt]\tiny $1/s$};

  \coordinate[right=8mm of int] (tap);
  \node[right=1mm of tap, lbl] (xist) {$x$};

  \draw[sig] ($(soll)+(-6mm,0)$) -- node[lbl,above,pos=0.3]{$x_{soll}$} (s);
  \draw[sig] (s)   -- (dz);
  \draw[sig] (dz)  -- (kp);
  \draw[sig] (kp)  -- (sat);
  \draw[sig] (sat) -- (rl);
  \draw[sig] (rl)  -- (tmc);
  \draw[sig] (tmc) -- (int);
  \draw[line] (int) -- (tap);
  \fill (tap) circle (1.2pt);

  \node[lbl, above left=-1mm and -2mm of s]     {$+$};
  \node[lbl, below left=-1.5mm and -3.5mm of s] {$-$};

  \node[msg, below=16mm of $(kp)!0.5!(sat)$] (enc)
        {Messglied -- AS5600\\[-1pt]\tiny 0,01304\,mm je Count, Median-3};

  \draw[sig] (tap) |- (enc.east);
  \draw[sig] (enc.west) -| (s);

  \node[grp] at ($(dz.north)!0.5!(sat.north)+(0,4.5mm)$)  {Regler und Grenzen};
  \node[grp] at ($(tmc.north)!0.5!(int.north)+(0,4.5mm)$) {Aktor und Strecke};

\end{tikzpicture}
\caption{Blockschaltbild des Positionsregelkreises. Der Integrator ist Teil der
         Regelstrecke, nicht des Reglers.}
\label{fig:positionsregelkreis}
\end{figure}
% ---------------------------------------------------------------

Tabelle \ref{tab:reglerparameter} stellt die Parameter des Reglers zusammen.

% ---------------------------------------------------------------
\begin{table}[htbp]
    \centering
    \caption{Parameter des Positionsreglers}
    \label{tab:reglerparameter}
    \begin{tabular}{llp{6cm}}
        \hline
        Größe & Wert & Herkunft \\
        \hline
        $T$        & 2~ms                    & Periode des Reglertasks \\
        $K_P$      & 125,7~s$^{-1}$          & $2\pi \cdot 20$~Hz \\
        $v_{max}$  & 40~mm/s                 & zulässige Verfahrgeschwindigkeit \\
        $a_{max}$  & 5000~mm/s²              & oberhalb $K_P^2 \cdot \Delta x$ \\
        Totzone    & 0,05~mm                 & rund vier Encoderstufen \\
        \hline
    \end{tabular}
\end{table}
% ---------------------------------------------------------------

% +++++++++++++++++++++++++++++++++++++++++++++++++++++++++++++++++
% Kraftregelung über StallGuard - Entwurf
% +++++++++++++++++++++++++++++++++++++++++++++++++++++++++++++++++
\section{Kraftregelung über StallGuard}
\label{sec:sg_entwurf}

Für die Greifkraft war ursprünglich ein zweiter, geschlossener Regelkreis vorgesehen. Da
die Aufgabenstellung die zunächst angedachte Strommessung ausschließt und ein Kraftsensor
nicht vorhanden ist, bot sich die im Treiber bereits implementierte StallGuard-Funktion
als Messglied an. Dieser Abschnitt beschreibt den Entwurf. Die messtechnische
Überprüfung und die daraus folgende Verwerfung des Konzepts werden in Abschnitt
\ref{sec:sg_verwerfung} behandelt.

\subsection*{SG\_RESULT als Messgröße}

Der Rotor eines Schrittmotors ist ein Permanentmagnet und induziert bei Drehung eine
Spannung in den Statorwicklungen. Unter Last bleibt er weiter hinter dem umlaufenden
Statorfeld zurück, wodurch sich die Phasenlage dieser Spannung verschiebt. Der TMC2209
wertet diesen Zusammenhang intern aus und stellt das Ergebnis als Registerwert
\texttt{SG\_RESULT} bereit.

Entscheidend für den Reglerentwurf ist das Vorzeichen dieser Kennlinie. StallGuard wurde
zur Erkennung eines blockierten Motors entwickelt, der ausgegebene Wert läuft daher mit
steigender Last gegen null. Ein hoher Messwert bedeutet also eine geringe Belastung. Für
die Fehlerbildung folgt daraus

\begin{equation}
    e = \mathrm{SG_{ist}} - \mathrm{SG_{soll}} ,
    \label{eq:sg_fehler}
\end{equation}

also Ist minus Soll und damit genau umgekehrt zum Positionsregler. Wird diese Umkehrung
übersehen, ergibt sich ein selbsterhaltender Stillstand: Der Regler berechnet eine
negative Stellgröße, die von der Begrenzung auf null abgeschnitten wird, der Antrieb
bewegt sich nicht, und die Messgröße ändert sich folglich ebenfalls nicht.

\subsection*{Die Wandlungskette}

Anders als beim Positionsregler, bei dem der Encoder die Regelgröße unmittelbar misst,
liegt zwischen Stellgröße und Messgröße eine Kette physikalischer Umwandlungen. Jedes
Glied dieser Kette ist ein eigener Modellblock:

\begin{align}
    \Delta x &= \max(x - x_K,\; 0) \label{eq:kette_weg}\\
    F        &= c_{ers} \cdot \Delta x \label{eq:kette_kraft}\\
    M        &= 2 \cdot r \cdot F + M_{reib} \label{eq:kette_moment}\\
    \mathrm{SG} &= \mathrm{SG_0} - k_{SG} \cdot M \label{eq:kette_sg}
\end{align}

Gleichung \ref{eq:kette_weg} bildet den Kontakt ab. Die Begrenzung auf nichtnegative Werte
ist nicht rechentechnischer Natur, sondern physikalisch: Ein nachgiebiges Objekt kann
drücken, aber nicht ziehen. Solange der Finger den Kontaktpunkt $x_K$ nicht erreicht hat,
ist die Kraft exakt null.

Der Faktor zwei in Gleichung \ref{eq:kette_moment} berücksichtigt, dass das Ritzel beide
Zahnstangen gleichzeitig antreibt und sich die Momente beider Finger am Motor addieren.
Bei einem Wälzkreisradius von $r = 8{,}5$~mm entspricht damit jedes Newton Greifkraft je
Finger einem Motormoment von 17~Nmm. Das Reibmoment $M_{reib}$ liegt auch bei Leerfahrt
an und wird deshalb additiv berücksichtigt und nicht in den Kontaktblock gezogen.

Für die Zielkraft von 5~N je Finger ergibt sich daraus ein Nutzmoment von 85~Nmm,
zuzüglich eines geschätzten Reibmoments von 15~Nmm also rund 100~Nmm. Bezogen auf das
Haltemoment des Motors von etwa 260~Nmm entspricht das einer Auslastung von 38~Prozent.

\subsection*{Reglerstruktur}

Der Regler selbst entspricht strukturell dem Positionsregler: Proportionalglied,
Stellgrößenbegrenzung, Anstiegsbegrenzung. Die Verstärkung $K_v$ ist in mm/s je Zählschritt
angegeben.

Ein Unterschied betrifft die Stellgrößenbegrenzung. Ihre untere Grenze liegt nicht bei
$-v_{max}$, sondern bei null. Der Greifer soll während des Kraftaufbaus unter keinen
Umständen zurückfahren, da ein Richtungswechsel zunächst das Flankenspiel der Verzahnung
durchläuft und die aufgebrachte Kraft dabei schlagartig auf null fällt. Der Kraftaufbau
ist damit eine ausschließlich monotone Zustellbewegung.

Für die Verstärkung selbst existiert keine Auslegungsformel analog zu Gleichung
\ref{eq:kp}, da die Strecke in diesem Zweig keine geschlossene Übertragungsfunktion
besitzt. $K_v$ wurde deshalb so gewählt, dass die Zufahrgeschwindigkeit bei unbelastetem
Antrieb oberhalb der Gültigkeitsgrenze von \texttt{SG\_RESULT} liegt: Unterhalb von etwa
5~mm/s ist die induzierte Spannung zu klein und der Messwert wird zu Rauschen.

\subsection*{Strukturelle Grenze des Konzepts}

\texttt{SG\_RESULT} wird nicht zeitgetaktet aktualisiert, sondern einmal je Vollschritt,
also alle 0,267~mm Verfahrweg. Aus dieser Eigenschaft folgen zwei Einschränkungen, die
bereits im Entwurf erkennbar waren.

Zum einen hängt die Zahl der Messwerte, die während des Kraftaufbaus überhaupt anfallen,
ausschließlich von der Nachgiebigkeit des Objekts ab:

\begin{equation}
    N = \frac{F_{soll}}{c_{ers} \cdot \Delta x_{voll}}
    \label{eq:messwerte}
\end{equation}

Die Geschwindigkeit kürzt sich heraus. Langsameres Zufahren erhöht die Zahl der Messwerte
also nicht. Fordert man mindestens zehn Stützstellen zwischen Kontakt und Zielkraft, damit
von einer Regelung und nicht von einem Schwellwertschalter gesprochen werden kann, folgt
daraus die Auslegungsgrenze

\begin{equation}
    c_{ers} < \frac{F_{soll}}{10 \cdot \Delta x_{voll}} = 1{,}87\,\mathrm{N/mm} .
    \label{eq:steifigkeitsgrenze}
\end{equation}

Oberhalb dieser Steifigkeit springt die gemessene Kraft von einem zu kleinen auf einen zu
großen Wert, ohne dass ein Zwischenschritt existiert, an dem der Regler eingreifen könnte.

Zum anderen öffnet sich der Kreis genau dann, wenn er sein Ziel erreicht. Mit
verschwindender Geschwindigkeit wächst das Abtastintervall über alle Grenzen, und die
Rückführung friert auf dem zuletzt gemessenen Wert ein. Was danach geschieht -- Kriechen
des Materials, Temperaturdrift, ein Lastwechsel -- ist strukturell unsichtbar. Formal
handelt es sich damit nicht um eine Kraftregelung im Arbeitspunkt, sondern um eine
gesteuerte Zustellbewegung mit ereignisgetriggertem Abbruch.

Beide Einschränkungen wurden in der Simulation bestätigt und in der Auslegung
berücksichtigt. Sie waren jedoch nicht der Grund für die spätere Verwerfung des Konzepts;
diese beruht auf einem Befund am Messglied selbst, der sich erst am Prüfstand zeigte
(Abschnitt \ref{sec:sg_verwerfung}).

% +++++++++++++++++++++++++++++++++++++++++++++++++++++++++++++++++
% Kraftregelung über Position und Federmodell
% +++++++++++++++++++++++++++++++++++++++++++++++++++++++++++++++++
\section{Kraftvorgabe über Position und Federmodell}
\label{sec:kraft_position}

Nachdem sich \texttt{SG\_RESULT} als Messglied für die auftretende Last als ungeeignet
erwiesen hat, wurde die Kraftaufgabe auf einem anderen Weg gelöst. Grundlage ist die
Überlegung, dass bei bekannter Nachgiebigkeit des Objekts Kraft und Weg über das
Hookesche Gesetz eindeutig zusammenhängen. Eine Kraftvorgabe lässt sich damit in eine
Positionsvorgabe umrechnen und an den bereits vorhandenen, validierten Positionsregler
übergeben.

\subsection*{Prinzip}

Ein eigener Kraftregelkreis existiert nicht. Die Zielposition ergibt sich aus dem
Kontaktpunkt $x_0$ und dem für die gewünschte Kraft erforderlichen Federweg:

\begin{equation}
    x_{soll} = x_0 + \frac{F_{soll}}{c_{eff}}
    \label{eq:kraft_vorsteuerung}
\end{equation}

Die tatsächlich aufgebrachte Kraft folgt umgekehrt aus der gemessenen Position:

\begin{equation}
    F = c_{eff} \cdot (x_{ist} - x_0)
    \label{eq:kraft_rueckrechnung}
\end{equation}

Bezogen auf die Kraft ist dies ein offener Kreis: Sie wird berechnet, nicht gemessen. Der
darin eingebettete Positionskreis bleibt dagegen geschlossen und korrigiert jede
Abweichung der Position. Abbildung \ref{fig:gesamtstruktur} zeigt diese Schachtelung.

% ---------------------------------------------------------------
\begin{figure}[htbp]
\centering
\begin{tikzpicture}[node distance=5mm, scale=0.92, transform shape]

  \node (fin) {};
  \node[ghost, right=9mm of fin] (gs)  {$\pm$};
  \node[blk, right=7mm of gs]    (inv) {$1/c_{e\!f\!f}$\\[-1pt]\tiny 2,12\,N/mm};
  \node[blk, right=5mm of inv]   (x0)  {$+\,x_0$\\[-1pt]\tiny Kontaktpunkt};
  \node[blkw, right=14mm of x0, fill=black!6] (pos)
        {\textbf{Positionsregelkreis}\\[-1pt]\tiny geschlossen, Abb.~\ref{fig:positionsregelkreis}};
  \node[blkw, right=14mm of pos] (fed)
        {Feder (Strecke)\\[-1pt]\tiny $F=c_{e\!f\!f}(x-x_0)$};

  \coordinate[right=8mm of fed] (fout);
  \node[right=1mm of fout, lbl] {$F$};

  \draw[sig] ($(fin)+(-6mm,0)$) -- node[lbl,above,pos=0.3]{$F_{soll}$} (gs);
  \draw[sig] (gs)  -- (inv);
  \draw[sig] (inv) -- (x0);
  \draw[sig] (x0)  -- node[lbl,above=0.5mm]{$x_{soll}$} (pos);
  \draw[sig] (pos) -- node[lbl,above=0.5mm]{$x_{ist}$} (fed);
  \draw[line] (fed) -- (fout);
  \fill (fout) circle (1.2pt);

  \begin{scope}[on background layer]
    \node[draw, dashed, black!45, rounded corners=2pt,
          inner sep=2mm, fit=(pos)] (innen) {};
  \end{scope}
  \node[grp, above=1mm of innen] {geregelt};

  \draw[line, dashed, black!55] (fout) -- ++(0,-16mm) coordinate (fb2);
  \draw[line, dashed, black!55] (fb2) -| ($(gs)+(0,-16mm)$);
  \draw[miss] ($(gs)+(0,-16mm)$) -- (gs);

  \coordinate (xm) at ($(fb2)!0.55!($(gs)+(0,-16mm)$)$);
  \fill[white] ($(xm)+(-2mm,-2mm)$) rectangle ($(xm)+(2mm,2mm)$);
  \draw[red!60!black, line width=0.9pt]
        ($(xm)+(-1.6mm,-1.6mm)$) -- ($(xm)+(1.6mm,1.6mm)$);
  \draw[red!60!black, line width=0.9pt]
        ($(xm)+(-1.6mm,1.6mm)$) -- ($(xm)+(1.6mm,-1.6mm)$);
  \node[font=\scriptsize, text=red!60!black, below=1mm of xm]
        {keine Kraftrückführung};

\end{tikzpicture}
\caption{Gesamtstruktur der Kraftvorgabe. Der Positionsregelkreis bildet den
         geschlossenen Kern; der Kraftpfad darum herum ist offen. Die gestrichelte
         Linie markiert die Stelle, an der eine echte Kraftregelung ihre Rückführung
         hätte.}
\label{fig:gesamtstruktur}
\end{figure}
% ---------------------------------------------------------------

\subsection*{Ablauf eines Greifvorgangs}

Der Kontaktpunkt $x_0$ ist der einzige Parameter, der nicht aus der Auslegung stammt,
sondern im Versuch bestimmt werden muss. Der Ablauf gliedert sich daher in vier Schritte:

\begin{enumerate}
    \item Der Finger wird positionsgeregelt zugestellt, bis das Objekt gerade berührt wird.
    \item Die aktuelle Position wird als Kontaktpunkt $x_0$ übernommen.
    \item Die Zielposition wird nach Gleichung \ref{eq:kraft_vorsteuerung} berechnet und
          angefahren.
    \item Der Antrieb hält diese Position. Die Kraft steht.
\end{enumerate}

Die Übernahme von $x_0$ erfolgt derzeit manuell über die Bedienoberfläche. Eine
automatische Erkennung ist möglich, indem der Lastwinkel -- die Differenz zwischen dem
Schrittzähler des Treibers und der gemessenen Encoderposition -- ausgewertet wird. Dieser
liegt bei Leerfahrt praktisch bei null und springt bei Berührung auf einen messbaren Wert.

Für die Zielkraft von 5~N ergibt sich bei der am Prüfstand wirksamen Steifigkeit von
$c_{eff} = 2{,}12$~N/mm ein Federweg von 2,36~mm, bei 10~N entsprechend 4,72~mm.

\subsection*{Genauigkeit und Fehlerempfindlichkeit}

Da die Kraft nicht gemessen wird, geht jeder Fehler im Modell oder in der Position
unmittelbar und unkorrigiert in die Kraft ein. Die erreichbare Genauigkeit folgt damit
direkt aus der Positioniergenauigkeit. Mit der Totzone von 0,05~mm aus Abschnitt
\ref{sec:positionsregler} ergibt sich

\begin{equation}
    \Delta F = c_{eff} \cdot \Delta x = 2{,}12\,\mathrm{N/mm} \cdot 0{,}05\,\mathrm{mm}
    = 0{,}11\,\mathrm{N} .
    \label{eq:kraftfehler}
\end{equation}

Bemerkenswert ist das Skalierungsverhalten: Die \emph{absolute} Kraftgenauigkeit
verschlechtert sich linear mit der Steifigkeit. Bei einer einzelnen Feder mit 1,06~N/mm
läge der Fehler bei nur 0,05~N. Konstant bleibt lediglich der auf den Federweg bezogene
\emph{relative} Fehler, da sich der erforderliche Federweg umgekehrt proportional zu
$c_{eff}$ verhält. Daraus folgt unmittelbar die Anwendungsgrenze des Verfahrens: An einem
sehr steifen Objekt wird der Federweg so klein, dass die Positionsauflösung die Kraft
nicht mehr sinnvoll auflöst.

Ein Vergleich der drei Fehlerquellen zeigt, welche Größe die Genauigkeit tatsächlich
bestimmt. Da alle Zusammenhänge linear sind, lassen sich die Beiträge unmittelbar
angeben; eine Parameterstudie ist dafür nicht erforderlich.

% ---------------------------------------------------------------
\begin{table}[htbp]
    \centering
    \caption{Einfluss der Fehlerquellen auf die Greifkraft bei $F_{soll} = 5$~N und
             $c_{eff} = 2{,}12$~N/mm}
    \label{tab:fehlerquellen}
    \begin{tabular}{lcc}
        \hline
        Fehlerquelle & Kraftfehler & bezogen auf $F_{soll}$ \\
        \hline
        Totzone, $\pm 0{,}05$~mm          & $\pm 0{,}11$~N & 2,1~\% \\
        Kontaktpunkt $x_0$ um 0,5~mm falsch & 1,06~N        & 21,2~\% \\
        $c_{eff}$ um 5~\% falsch           & 0,25~N        & 5,0~\% \\
        \hline
    \end{tabular}
\end{table}
% ---------------------------------------------------------------

Die Bestimmung des Kontaktpunkts ist damit die kritische Größe des Verfahrens. Ein Fehler
von einem halben Millimeter, wie er bei manueller Zustellung ohne weitere Hilfsmittel
durchaus auftreten kann, wirkt sich zehnmal stärker aus als die gesamte
Positionierungenauigkeit des Reglers. Die Güte der Kraftvorgabe hängt folglich weniger am
Regler als an der Reproduzierbarkeit des Kontaktpunkts und an der Kenntnis der
Steifigkeit.

\subsection*{Einordnung}

Formal handelt es sich bei dem beschriebenen Verfahren um eine \emph{modellbasierte
Vorsteuerung mit geregelter Position} und nicht um eine Kraftregelung. Die Unterscheidung
ist nicht nur begrifflich, sondern bestimmt, welche Störungen ausgeregelt werden und
welche nicht:

\begin{itemize}
    \item \textbf{Keine Rückführung der Kraft.} Alles, was nach dem Anfahren der
          Zielposition geschieht -- Kriechen des Materials, Temperaturdrift, ein
          Lastwechsel am Objekt -- bleibt unbemerkt.
    \item \textbf{Die Steifigkeit muss bekannt sein.} An einem unbekannten Objekt ist das
          Verfahren nicht anwendbar, ohne $c_{eff}$ vorher zu bestimmen.
    \item \textbf{Der Kontaktpunkt geht direkt in die Kraft ein.} Er ist kein
          Auslegungsparameter, sondern muss bei jedem Greifvorgang neu ermittelt werden.
\end{itemize}

Diesen Einschränkungen steht gegenüber, dass das Verfahren ohne zusätzliche Sensorik
auskommt, auf einem bereits validierten Regelkreis aufsetzt und die geforderte
Umschaltung zwischen Positions- und Kraftbetrieb strukturell trivial macht: Beide
Betriebsarten nutzen denselben Regler und unterscheiden sich lediglich in der Herkunft
des Sollwerts. Die messtechnische Überprüfung erfolgt in Abschnitt
\ref{sec:ergebnis_kraft}.

% +++++++++++++++++++++++++++++++++++++++++++++++++++++++++++++++++
% Umschaltlogik und Zustandsautomat
% +++++++++++++++++++++++++++++++++++++++++++++++++++++++++++++++++
\section{Umschaltlogik und Zustandsautomat}
\label{sec:zustandsautomat}

Der Regler allein bildet nur einen Ausschnitt eines Greifvorgangs ab. Er berechnet aus
einer Abweichung eine Geschwindigkeit, weiß aber nicht, woher sein Sollwert stammt, ob
gerade referenziert, angenähert oder gegriffen wird und wie im Fehlerfall zu reagieren
ist. Diese Entscheidungen trifft eine übergeordnete Ablaufsteuerung, die als
Zustandsautomat im selben 2-ms-Task ausgeführt wird.

\subsection*{Zustände}

Die Ablaufsteuerung umfasst sechs Zustände. Tabelle \ref{tab:zustaende} fasst ihre
Aufgaben zusammen, Abbildung \ref{fig:zustandsautomat} zeigt die Übergänge.

% ---------------------------------------------------------------
\begin{table}[htbp]
    \centering
    \caption{Zustände der Ablaufsteuerung}
    \label{tab:zustaende}
    \begin{tabular}{llp{7.5cm}}
        \hline
        Zustand & Nr. & Aufgabe \\
        \hline
        \texttt{ST\_ZERO}      & 1 & Referenzfahrt, Nullpunkt setzen \\
        \texttt{ST\_DIRDETECT} & 2 & Bestimmung der Encoderrichtung \\
        \texttt{ST\_CONTROL}   & 3 & regulärer Regelbetrieb, Position oder Kraft \\
        \texttt{ST\_HOLD}      & 4 & Halten der erreichten Position bzw. Kraft \\
        \texttt{ST\_FAULT}     & 5 & Fehlerzustand, Motor gestoppt \\
        \texttt{ST\_OPEN}      & 6 & Öffnen und Lösen des Objekts \\
        \hline
    \end{tabular}
\end{table}
% ---------------------------------------------------------------

% ---------------------------------------------------------------
\begin{figure}[htbp]
\centering
\begin{tikzpicture}[x=1cm, y=1cm]

  \node[zst, fill=white] (idle) at (0,0)   {Idle\\[-1pt]\tiny nach dem Start};
  \node[zst]  (zero) at (3.4,0)            {\texttt{ST\_ZERO}\\[-1pt]\tiny Referenzfahrt};
  \node[zst]  (dir)  at (6.8,0)            {\texttt{ST\_DIRDETECT}\\[-1pt]\tiny Richtungs\-erkennung};
  \node[zst]  (ctrl) at (10.2,0)           {\texttt{ST\_CONTROL}\\[-1pt]\tiny Regelbetrieb};

  \node[zflt] (flt)  at (1.7,-2.4)         {\texttt{ST\_FAULT}\\[-1pt]\tiny Motor gestoppt};
  \node[zst]  (open) at (6.8,-2.4)         {\texttt{ST\_OPEN}\\[-1pt]\tiny Öffnen};
  \node[zst]  (hold) at (10.2,-2.4)        {\texttt{ST\_HOLD}\\[-1pt]\tiny Kraft halten};

  \draw[ztr] (idle) -- node[above]{Kommando}      (zero);
  \draw[ztr] (zero) -- node[above]{Nullpunkt}     (dir);
  \draw[ztr] (dir)  -- node[above]{Richtung}      (ctrl);
  \draw[ztr] (ctrl) -- node[right]{Ziel erreicht} (hold);
  \draw[ztr] (hold) -- node[below]{Kommando}      (open);
  \draw[ztr] (open.north east) -- node[above,pos=0.55,sloped]{fertig} (ctrl.south west);

  \draw[ztrd] (zero.south)     -- (flt.north east);
  \draw[ztrd] (dir.south west) -- (flt.north);
  \draw[ztrd] (open.west)      -- (flt.east);

  \node[grp, below=3mm of flt, text width=54mm, align=left]
        {gestrichelt: Fehlerbedingung. Sie führt aus \emph{jedem}
         Zustand nach \texttt{ST\_FAULT}; verlassen wird dieser nur
         durch einen Neustart.};

\end{tikzpicture}
\caption{Zustände und Übergänge der Ablaufsteuerung.}
\label{fig:zustandsautomat}
\end{figure}
% ---------------------------------------------------------------

\subsection*{Bestimmung der Encoderrichtung}

Der Zustand \texttt{ST\_DIRDETECT} verdient besondere Erwähnung, da er im ursprünglichen
Entwurf nicht vorgesehen war. Die Zuordnung zwischen der Zählrichtung des Encoders und
der Verfahrrichtung des Antriebs hängt davon ab, wie der Magnet auf der Welle sitzt und
wie der Sensor gegenüber der Mechanik orientiert ist. Sie ist aus der Konstruktion nicht
zuverlässig ableitbar.

Diese Zuordnung ist kein Randdetail. Bei falschem Vorzeichen wird aus der Gegenkopplung
des Regelkreises eine Mitkopplung: Die Regelabweichung wächst mit jedem Takt, der Regler
erhöht daraufhin die Geschwindigkeit weiter, und der Antrieb läuft weg, statt zu regeln.
Statt die Richtung als Konstante im Code festzulegen und bei einem Umbau auf einen Fehler
zu warten, wird sie beim Start experimentell bestimmt: Der Antrieb fährt mit einer kleinen,
festen Geschwindigkeit, und aus dem Vorzeichen der Encoderänderung wird die Zuordnung
abgeleitet.

\subsection*{Umschaltung zwischen den Betriebsarten}

Die Anforderung einer nahtlosen Umschaltung zwischen Positions- und Kraftbetrieb während
des laufenden Betriebs wird über die Variable \texttt{g\_ctrl\_mode} umgesetzt, die zur
Laufzeit über ein Bedienkommando geändert werden kann. Beide Betriebsarten laufen im
Zustand \texttt{ST\_CONTROL}.

Durch die in Abschnitt \ref{sec:kraft_position} gewählte Lösung ist diese Umschaltung
strukturell trivial. Da die Kraftvorgabe keinen eigenen Regelkreis besitzt, sondern
lediglich den Sollwert des Positionsreglers auf andere Weise berechnet, wird beim
Umschalten kein Regler getauscht. Es ändert sich ausschließlich die Herkunft von
$x_{soll}$: im Positionsbetrieb ein direkt vorgegebener Wert, im Kraftbetrieb das Ergebnis
von Gleichung \ref{eq:kraft_vorsteuerung}.

Damit entfallen sämtliche Probleme, die bei der Umschaltung zwischen zwei getrennten
Reglern zu behandeln wären: Es gibt keine Zustandsgrößen zu übergeben, keinen
stoßfreien Übergang zu konstruieren und keinen Sprung der Stellgröße im Umschaltmoment.
Die ursprünglich mit \texttt{SG\_RESULT} geplante Lösung hätte diesen Aufwand erfordert,
da dort ein zweiter Regler mit eigener Rückführung und eigener Fehlerbildung aktiv gewesen
wäre. Der Wegfall des zweiten Regelkreises ist an dieser Stelle kein Verlust, sondern eine
Vereinfachung.

\subsection*{Halten der Greifkraft}

Nach Erreichen der Zielposition geht die Ablaufsteuerung in den Zustand
\texttt{ST\_HOLD} über. Der Antrieb steht, muss aber weiterhin ein Moment aufbringen, da
die gespannte Feder ihn zurückzudrücken versucht.

Motortreiber reduzieren den Spulenstrom im Stillstand üblicherweise auf einen niedrigeren
Haltestrom, um die thermische Belastung zu senken. Der TMC2209 schaltet dazu nach Ablauf
der Zeit \texttt{TPOWERDOWN} von \texttt{IRUN} auf \texttt{IHOLD} um. Bei der hier
vorliegenden Anwendung führt dieses Verhalten jedoch dazu, dass das verbleibende
Haltemoment die anliegende Federkraft nicht mehr sicher aufnimmt. Die Folge wäre ein
Zurückdrücken des Antriebs und damit ein Einbruch der Greifkraft.

Aus diesem Grund wurde \texttt{IHOLD} dauerhaft auf denselben Wert wie \texttt{IRUN}
gesetzt. Gegenüber der ursprünglich vorgesehenen Umschaltung zur Laufzeit ist diese
Lösung einfacher und vermeidet einen weiteren Nebeneffekt: Die Kennlinie von
\texttt{SG\_RESULT} ist stromabhängig, sodass ein Wechsel des Stroms während des Betriebs
die Messwerte auf unterschiedliche Kennlinien legen würde. Erkauft wird dies mit einer
höheren Dauerverlustleistung im Stillstand, die bei den vorliegenden Strömen unkritisch
ist.

\subsection*{Fehlerverhalten}

Aus jedem Zustand führt eine Fehlerbedingung nach \texttt{ST\_FAULT}. Dort wird der Motor
angehalten, und weder Positions- noch Kraftvorgabe greifen weiterhin auf die
Motoransteuerung zu. Der Zustand kann nicht über ein reguläres Bedienkommando verlassen
werden; für eine erneute Inbetriebnahme ist ein bewusster Neustart erforderlich, nach dem
Referenzfahrt und Richtungserkennung erneut durchlaufen werden. Die einzelnen
Fehlerbedingungen und ihre Erkennung sind in Abschnitt \ref{sec:ueberlastschutz}
beschrieben.
